\chapter{Paarwerte in Mengenfamilien}
\label{b14guide:ch:sets}

\section{Paarwerte relativ zu einer Familie}

\begin{definition}[Paaroperationen einer Mengenfamilie]
\label{b14guide:def:family}
Sei $U$ eine Menge und $\mathcal M\subseteq\mathcal P(U)$
eine beliebige Familie, versehen mit der Inklusionsordnung.
Für $X,Y\in\mathcal M$ bezeichnen
$X\wedge_{\mathcal M}Y$ und $X\vee_{\mathcal M}Y$
die Paarwerte im Träger $\mathcal M$, sofern sie existieren.
Die Familie wird nicht als abgeschlossen unter Durchschnitt
oder Vereinigung vorausgesetzt.
\end{definition}

\begin{theorem}[Größte enthaltene und kleinste enthaltende Familienmenge]
\label{b14guide:thm:family-characterization}
Für $I,S\in\mathcal M$ gilt
\[
\begin{aligned}
 I=X\wedge_{\mathcal M}Y
 &\Longleftrightarrow
 \bigl(I\subseteq X\cap Y\ \land\
 \forall Z\in\mathcal M\;(Z\subseteq X\cap Y\Rightarrow Z\subseteq I)\bigr),\\
 S=X\vee_{\mathcal M}Y
 &\Longleftrightarrow
 \bigl(X\cup Y\subseteq S\ \land\
 \forall Z\in\mathcal M\;(X\cup Y\subseteq Z\Rightarrow S\subseteq Z)\bigr).
\end{aligned}
\]
Ein vorhandener Paarwert kann deshalb vom mengentheoretischen
Durchschnitt beziehungsweise von der Vereinigung abweichen.
\end{theorem}

\begin{theorem}[Durchschnitt und Vereinigung als Paarwerte]
\label{b14guide:thm:family-operations}
Für $X,Y\in\mathcal M$ gilt
\[
\begin{aligned}
 X\cap Y\in\mathcal M
 &\Longrightarrow X\wedge_{\mathcal M}Y=X\cap Y,\\
 X\cup Y\in\mathcal M
 &\Longrightarrow X\vee_{\mathcal M}Y=X\cup Y.
\end{aligned}
\]
Jede Implikation sichert die Existenz des betreffenden Paarwertes.
Ein vorhandener Wert ist genau dann der mengentheoretische
Durchschnitt beziehungsweise die Vereinigung, wenn diese Menge
selbst zur Familie gehört.
\end{theorem}

\begin{theorem}[Paarwerte im Potenzmengenträger]
\label{b14guide:thm:powerset}
Für alle $X,Y\subseteq U$ existieren im Träger $\mathcal P(U)$
beide Paarwerte, und es gilt
\[
 X\wedge_{\mathcal P(U)}Y=X\cap Y,\qquad
 X\vee_{\mathcal P(U)}Y=X\cup Y.
\]
Diese Identifikation gilt für jedes einzelne Paar, ohne daß
für $\mathcal P(U)$ eine weitere axiomatische Theorie
eingeführt werden muß.
\end{theorem}